\section{Fourier-feature coordinate charts}\label{sec:fourier}

The transition map of Section~\ref{sec:tmap} reads contact off the boundary charts $\map{A}$ and
$\map{B}$. Its accuracy is therefore the accuracy of those charts as representations of the geometry. A
chart that smooths the asperity slopes returns a smoothed gap and a smoothed normal, and no
constraint-enforcement scheme downstream can recover the lost detail. This section states how the charts
are built and trained, and isolates the one ingredient that lets a small network resolve the sharp, rough,
cusped boundaries that an ambient level set (and a plain coordinate network) smooth away: a fixed
Fourier-feature encoding of the input. The level set, a neural signed-distance function, appears here only
as the baseline that the encoding has to beat. We treat it as a fair baseline, not a strawman, because it is
itself a coordinate network and suffers the same pathology a plain chart does.

\subsection{Spectral bias of a coordinate network}\label{sec:spectral}

A coordinate network is a function $f_\theta:\R^d\!\to\!\R$ given by a stack of affine maps and $\tanh$
nonlinearities fed the raw coordinate. Two of the objects this paper compares are exactly of this form:
an ambient neural signed-distance function $\Phi_\theta(\xbf)$ of the spatial point $\xbf\in\R^3$, and a
boundary chart $h_\theta(x)$ of the surface parameter $x$. Both are trained by regression onto sampled
data, and both inherit the same defect.

% code: rahaman2019spectral
A coordinate network learns the low-frequency content of its target first and the high-frequency content
slowly or not at all. This is the spectral bias of \citet{rahaman2019spectral}: in the
neural-tangent-kernel description of wide-network training, the kernel of a coordinate network is
approximately translation invariant and Laplace-like, so its eigenvalues decay with frequency and the
error in each Fourier mode of the target decays at a rate set by the corresponding eigenvalue. We can
make this rate explicit. Write the target as $\hat z(k)$ in frequency $k$ and the network output after
training time $t$ as $\hat f_\theta(k,t)$; the residual in each mode then shrinks geometrically, at its
own frequency-dependent speed,
%
% code: rahaman2019spectral
\begin{equation}\label{eq:ntk-decay}
\hat f_\theta(k,t)-\hat z(k)\ \approx\ \big(\hat f_\theta(k,0)-\hat z(k)\big)\,e^{-\eta\,\lambda(k)\,t},
\qquad \lambda(k)\sim \norm{k}^{-(d+1)},
\end{equation}
%
where $\eta$ is the learning rate and $\lambda(k)$ is the kernel eigenvalue for frequency $k$. The moving
part is $\lambda(k)$: it sets how fast mode $k$ converges, and because it falls off as a power of
$\norm{k}$, the high-$\norm{k}$ modes converge at a vanishing rate. At any finite training budget the
recovered surface is therefore a low-pass-filtered version of the true one. On a rough surface the
consequence is geometric. The asperities, cusps, and fine relief are filtered out, and what survives
is a smoothed envelope.

This is the failure mode of the ambient level set, and it is intrinsic rather than incidental. A neural SDF
cannot represent a boundary sharper than the highest frequency its training has resolved through
\eqref{eq:ntk-decay}; on a rock joint it flattens the asperity flanks, and through the strength law
$\muapp=\tan(\phi_b+i)$ of \citet{patton1966multiple} a flattened asperity angle $i$ directly
under-predicts the friction. The point to make precise, and the basis for the controlled ablation below,
is that \emph{a boundary chart of the raw coordinate inherits the identical pathology}. Representing the
geometry as a function on the surface's own parameter, rather than as an ambient field, does not by itself
buy any high frequencies; it changes the domain of $f_\theta$, not its spectral bias. The chart
parametrization alone does not beat the level set. The Fourier features do.

\subsection{The Fourier-feature encoding and three banks}\label{sec:banks}

% code: tancik2020fourier
The remedy of \citet{tancik2020fourier} is to replace the raw coordinate with a fixed (non-trainable) bank
of sinusoids applied before the network. Let $\tilde x\in[-1,1]^d$ be the input normalized to the unit box,
%
\begin{equation}\label{eq:normalize}
\tilde x = \frac{2(x-x_{\mathrm{lo}})}{x_{\mathrm{hi}}-x_{\mathrm{lo}}}-1,
\end{equation}
%
and let $B\in\R^{K\times d}$ be a bank of $K$ frequency vectors. The encoding and the chart are
%
% code: tancik2020fourier
\begin{equation}\label{eq:gamma}
\gamma(x)=\big[\cos(2\pi B\tilde x),\ \sin(2\pi B\tilde x)\big]\in\R^{2K},
\qquad
h_\theta(x)=\mathtt{base}+\mathrm{MLP}_\theta\!\big(\gamma(x)\big).
\end{equation}
%
The bank $B$ is a registered buffer, fixed for the life of the chart; only the network weights $\theta$ and
the scalar offset $\mathtt{base}$ (which carries the mean height, so the network learns only the deviation)
are trained. Prepending $\gamma$ changes the induced kernel from the Laplace-like form of
\eqref{eq:ntk-decay} to a stationary, band-limited one whose bandwidth is set by $B$. To see the new
kernel concretely, take the one-dimensional geometric bank below: the kernel becomes a finite sum of
cosines, one per chosen frequency $f_k$, and it depends only on the separation $\tilde x-\tilde x'$,
%
% code: tancik2020fourier
\begin{equation}\label{eq:kernel}
\Theta(x,x')=\frac1K\sum_{k=1}^{K}\cos\!\big(2\pi f_k\,(\tilde x-\tilde x')\big),
\end{equation}
%
so the network's similarity measure is built from exactly the frequencies in the bank and no others. For
the isotropic Gaussian bank the kernel is instead the radial profile $J_0(2\pi\sigma\norm{\tilde{\bm u}-\tilde{\bm
v}})$, flat out to a cutoff $\sim 2\pi\sigma$. Every frequency below the cutoff is now learned in parallel,
at a comparable rate, instead of in slow frequency order. The network no longer has to discover the high
frequencies; they are handed to it, and it learns only their amplitudes and phases. Two properties make
this the right tool for geometry. First, the cutoff is transparent and tunable. Where the level set's
smoothing is an intrinsic property of its architecture, the Fourier cutoff is an explicit knob ($f_{\max}$,
$\sigma$, or the harmonic order $K$), set to just resolve the finest wavelength required. Second, the bank
is deterministic and frozen, so the encoding is a reproducible preprocessing of the input rather than a
learned, run-dependent representation.

% code: solvers/contact/profile_chart_2d.py::NeuralHeight1D ; solvers/contact/surface_chart_3d.py::NeuralHeight2D ; solvers/contact/radial_chart_2d.py::NeuralRho2D
This paper uses three charts, and they differ only in two respects: the domain each one lives on, and
which bank of frequencies it uses to tile frequency space. The training loop is identical across all
three. Each bank is matched to its geometry. The geometric bank gives equal coverage on a
\emph{log}-frequency axis, which suits a profile with asperities spread across many wavelengths; the
Gaussian bank gives isotropic two-dimensional coverage, which suits a height field over a patch; and the
harmonic bank is periodic by construction, so a chart wrapped around the circle closes up on itself.
Table~\ref{tab:banks} collects the three banks, the charts they encode, and their parameters.

\begin{table}
\centering
\caption{The three Fourier banks and the charts they encode. The geometric bank gives equal coverage of
\emph{log}-scale, the Gaussian bank gives isotropic two-dimensional coverage, and the harmonic bank is
periodic by construction so a chart on the circle closes up. The bank $B$ is fixed; only the network and a
scalar offset are trained.}
\label{tab:banks}
\begin{tabular}{@{}llll@{}}
\toprule
Bank & Definition & Chart (file::class) & Parameters \\
\midrule
Geometric (deterministic) & $B=\operatorname{geomspace}(0.5,f_{\max},K)$ &
\code{NeuralHeight1D} ($h(x)$, $x\in\R$) & $K{=}192$, $f_{\max}{=}1500$ \\
Gaussian random & $B_i\sim\mathcal N(0,\sigma^2 I)$ &
\code{NeuralHeight2D} ($h(x,y)$) & $K{=}256$, $\sigma{=}8$ \\
Harmonic (periodic) & $[\cos(k\psi),\sin(k\psi)]_{k=1}^{K}$ &
\code{NeuralRho2D} ($\rho(\psi)$, $\psi\in\Sph^1$) & $K{=}16$ \\
\bottomrule
\end{tabular}
\end{table}

% code: solvers/contact/profile_chart_2d.py::NeuralHeight1D
The \emph{geometric} bank, $B=\operatorname{geomspace}(0.5,f_{\max},K)$ with $K=192$ and $f_{\max}=1500$,
spaces its frequencies equally in $\log$-scale and so resolves a wide range of asperity wavelengths with a
modest feature count; it carries the one-dimensional height chart $h_\theta(x)$ used for the rock-joint
profile, where the measured surface requires on the order of $L/\Delta x\approx 3000$ resolvable
wavelengths.
% code: solvers/contact/surface_chart_3d.py::NeuralHeight2D
The \emph{Gaussian random} bank, $B_i\sim\mathcal N(0,\sigma^2 I)$ with $K=256$ and $\sigma=8$, gives the
isotropic two-dimensional coverage appropriate to a height field $h_\theta(x,y)$ over a patch; its induced
kernel $J_0(2\pi\sigma\norm{\cdot})$ is flat to $\sim 2\pi\sigma$, the same Gaussian-relief encoding used by
the boundary-fitted volumetric decoder of Section~\ref{sec:form}.
% code: solvers/contact/radial_chart_2d.py::NeuralRho2D
The \emph{harmonic} bank, $[\cos(k\psi),\sin(k\psi)]_{k=1}^{K}$ with $K=16$, is periodic by construction;
this is mandatory, not a convenience, for the radial chart $\rho_\theta(\psi)$ on the circle
$\Sph^1$, where the radius must close up at $\psi=0,2\pi$. The radial chart additionally wraps its output in
a softplus, $\rho=\mathrm{softplus}(\mathtt{base}+\Delta)$, to keep $\rho>0$. Figure~\ref{fig:arch} shows
the encoding and the network it feeds.

\begin{figure}
\centering
\resizebox{\linewidth}{!}{%
\begin{tikzpicture}[font=\footnotesize,>={Stealth[length=2.2mm]},
  box/.style={draw,rounded corners=2pt,align=center,inner sep=4pt,minimum height=12mm,line width=0.7pt},
  tag/.style={draw,rounded corners=2pt,align=center,inner sep=3pt,font=\scriptsize}]
  \node[box,fill=gray!8] (in) {input\\ $x$};
  \node[box,fill=gray!8,right=5mm of in] (nm) {normalize\\ $\tilde x\in[-1,1]$};
  \node[box,draw=teal!70!black,fill=teal!6,right=5mm of nm] (ff)
    {Fourier features\\ $\gamma(x)=[\cos 2\pi B\tilde x,\ \sin 2\pi B\tilde x]$};
  \node[box,draw=blue!70!black,fill=blue!6,right=5mm of ff] (mlp)
    {$\mathrm{MLP}_\theta$\\ tanh, width 128, depth 4};
  \node[box,draw=teal!70!black,fill=teal!6,right=5mm of mlp] (out)
    {$h_\theta(x)=\mathrm{base}+\Delta$};
  \draw[->] (in) -- (nm);
  \draw[->] (nm) -- (ff);
  \draw[->] (ff) -- (mlp);
  \draw[->] (mlp) -- (out) node[midway,above,font=\scriptsize]{$+\,\mathrm{base}$};
  \node[tag,draw=orange!80!black,fill=orange!8,below=5mm of ff] (t1) {FIXED: bank $B$ not trained};
  \node[tag,draw=blue!70!black,fill=blue!6,below=5mm of mlp,xshift=8mm] (t2)
    {TRAINED: weights $\theta$ + scalar base (Adam, cosine, MSE)};
  \draw[orange!80!black,->,dashed] (t1) -- (ff);
  \draw[blue!70!black,->,dashed] (t2) -- (mlp);
\end{tikzpicture}%
}
\caption{The Fourier-feature chart architecture, drawn for the one-dimensional height chart
\code{NeuralHeight1D}. The coordinate is normalized to $[-1,1]$ via \eqref{eq:normalize} and expanded by
the \emph{fixed} bank \eqref{eq:gamma} into $2K=384$ cosine/sine features; a plain $\tanh$ network (depth
$4$, width $128$) maps those to a scalar deviation $\Delta$, offset by the learnable mean $\mathtt{base}$.
Only the network weights $\theta$ and $\mathtt{base}$ are trained; the bank $B$ is frozen. The
two-dimensional and three-dimensional charts replace the geometric bank with a Gaussian one and the radial
chart wraps the output in a softplus, leaving the rest unchanged.}
\label{fig:arch}
\end{figure}

\subsection{Training and the verify-first protocol}\label{sec:train}

Every chart is trained by the same procedure: normalize the input by \eqref{eq:normalize}, encode with the
fixed bank \eqref{eq:gamma}, then regress the network onto sampled surface points by minibatch
mean-squared error,
%
\begin{equation}\label{eq:mse}
\mathcal L(\theta)=\frac{1}{\abs{\mathcal B}}\sum_{i\in\mathcal B}\big(h_\theta(x_i)-z_i\big)^2,
\end{equation}
%
with Adam at learning rate $2\times10^{-3}$ and a cosine annealing schedule that decays the rate to zero
over the iteration budget. The networks are small and are trained in double precision on CPU, which
conditions the geometry and the Jacobians used downstream; the random seed is fixed so a run is
reproducible. There is no adversary, no loss schedule, and no curriculum: a plain MSE regression suffices
because the Fourier encoding, not the optimizer, carries the high frequencies.

A trained chart \emph{is} geometry, and a corrupted chart silently corrupts every downstream solve, so a
chart is admitted to a contact computation only after three checks. (i) \emph{Reconstruction fidelity}: the
RMSE on the sampled surface points, reported as a fraction of the surface RMS, must reach a few percent.
(ii) \emph{No element foldover}: for the volumetric decoder the chart-Jacobian determinant must satisfy
$\det J>0$ on every element, since a folded chart has no valid pushforward. (iii) \emph{Discretization
convergence}: a manufactured-solution study on the curved chart geometry must recover the $O(h^2)$ rate of
the underlying P1 chart finite-element method, confirming that the trained geometry did not spoil the
solver. A chart that fails, with relief amplitude too high so $\det J$ folds, or cutoff too low so the
surface is under-resolved, is rejected, not patched.

% code: solvers/contact/profile_chart_2d.py::NeuralHeight1D
The clearest measurement of the Fourier contribution is the framework's own ablation, run on the real
Inada-granite tensile-fracture profile \citep{inada2020digitalrocks}, a self-affine surface with RMS
$\approx 1.7$~mm. The same network, with identical width, depth, optimizer, and schedule, is trained
twice on the measured points, once with the Fourier features active and once on the raw coordinate (the
\code{plain=True} control). Because a measured profile is a deterministic surface to represent, the metric
is reconstruction RMSE on the measured points, and the bank is band-limited below the sampling Nyquist so
the fit is well posed rather than an interpolation to zero. Figure~\ref{fig:curves} reports the two
training curves and the recovered surfaces. With the Fourier features the chart descends through the
asperity scales to a reconstruction RMSE of $0.003\%$ of the RMS and recovers this ablation window's mean
slope angle, here $\bar i = 19.9^\circ$ (the full-profile capstone below measures the surface's mean angle
as $19.4^\circ$). The byte-for-byte identical network on the raw coordinate plateaus at a spectral-bias
floor of $7.6\%$ of RMS ($136~\mu\mathrm{m}$) and
low-passes the slopes to $\bar i = 11.5^\circ$. The two networks differ only in the input encoding, and one
converges to the surface while the other stalls at a smoothed approximation. This is the chart-versus-level-set
thesis reduced to its mechanism, on a real surface.

\begin{figure}
\centering
\includegraphics[width=\linewidth]{fourier_training_curves_pub}
\caption{Genuine training on the real Inada-granite profile \citep{inada2020digitalrocks}, both runs using
the repository's own \code{NeuralHeight1D}. (a) Reconstruction RMSE versus iteration: the Fourier chart
descends to $0.003\%$ of the surface RMS, while the \emph{identical} network fed the raw coordinate
plateaus at a $7.6\%$ ($136~\mu$m) spectral-bias floor; only the input encoding differs. (b) The learned
surfaces: the Fourier chart tracks the asperities and recovers the profile's mean slope angle $\bar i=19.9^\circ$,
whereas the plain network low-passes them to $\bar i=11.5^\circ$, under-resolving the slopes that set the
friction angle through $\muapp=\tan(\phi_b+i)$.}
\label{fig:curves}
\end{figure}

% code: solvers/fem/rough_block_decoder.py
The same separation holds for the boundary-fitted volumetric decoder used by the chart finite-element
method of Section~\ref{sec:form}, where the relief is carried by the Gaussian bank. The Fourier
decoder reconstructs the rough face to $2.2\%$ of RMS, a vanilla $\tanh$ decoder of equal capacity reaches
only $48\%$ (the spectral-bias floor again), and an ambient three-dimensional neural SDF of the same surface
reaches $15.7\%$; the decoder additionally passes $\det J>0$ on every element and the $O(h^2)$
manufactured-solution check (measured rates $2.07$ and $1.83$). On the one-dimensional Inada profile, the
production height chart ($K=192$, $f_{\max}=1500$) reconstructs the surface to $2.3~\mu\mathrm{m}$, about
$0.1\%$ of the $1.7$~mm RMS, while the ambient neural SDF of the same surface reaches only
$107~\mu\mathrm{m}$, which is $47\times$ worse with \emph{more} parameters, and smooths the mean asperity
angle from $19.4^\circ$ to $12.5^\circ$. Through the Patton law that smoothing under-predicts the peak shear
strength by $61\%$; the asperity slopes the chart preserves are precisely the geometry that drives the
contact mechanics.

\subsection{Why charts beat the level set}\label{sec:why}

The advantage is a property of the kernel, not of the parametrization. The level set is a coordinate
network of the ambient point and therefore carries the Laplace-like kernel of \eqref{eq:ntk-decay}, whose
spectral ceiling is an intrinsic, non-tunable consequence of the architecture; on rough or cusped geometry
it filters away exactly the high frequencies that contact cares about. The Fourier-feature chart carries
the stationary, band-limited kernel of \eqref{eq:kernel} instead, with a cutoff one sets explicitly to
resolve the finest wavelength of interest and no spectral ceiling below it. The chart also stores geometry
as a surface function, so its cost scales with the surface rather than with an ambient volume, but the
fidelity argument rests on the kernel.

The ablation of Section~\ref{sec:train} is the load-bearing evidence: the Fourier features, not the chart
parametrization alone, do the work. A chart of the raw coordinate (\code{plain=True}) smooths the Inada
profile to the same $7.6\%$ floor and the same $11.5^\circ$ angle that the level set produces, because it
inherits the same spectral bias; only when the fixed sinusoidal bank is prepended does the reconstruction
fall to $0.003\%$ and the slope angle become exact. We therefore do not claim that reading geometry off a
boundary chart is, by itself, more accurate than reading it off a level set. On a smooth surface dense
interpolation is easy for either, and the height RMSE of a well-resolved SDF is not the discriminating
quantity. The claim is sharper and is the one the measurements support: the level set smooths asperity
slopes through an intrinsic spectral ceiling and so under-predicts strength and dilatancy, whereas a
Fourier-feature chart has a transparent, tunable cutoff and resolves those slopes. The scope is that
the height and surface charts represent single-valued geometry (no overhangs) and that the harmonic radial
chart requires star-shaped boundaries; the band-limit is a chosen resolution, not an unlimited one. Within
that scope the encoding is what converts a smoothing approximator into a faithful one, and every downstream
result in this paper rests on it.
